\subsection{最初の非零共有chartの全必要係数被覆}
\label{subsec:cubic-CY-four-point-first-cover}
\begin{theorem}\label{thm:cubic-CY-four-point-first-empty}
四単純基点の元係数表示で $x_1\ne0,x_0=x_3=x_4=0$ を満たす領域と、
その全monomial軌道はH3包絡を与えない。
\end{theorem}
\begin{proof}
全支持作用は補題\ref{lem:cubic-CY-four-point-support}で確認した。
$x_1=1$ とし、残る二つの独立な対角比で
$(x_2,x_5)=(0,0),(0,1),(1,0),(1,1)$ の四枝へ移す。
以下の被覆は固定出典\cite{AS4CYpinned}のFN1--FN6を用い、
全21係数をあらかじめ零にせず扱う。
当repoでも元Schur行列から四枝の全逐次消去を再生した。
各消去で記録した式が実際の残余式であること、一次の主係数が
有理定数と $\lambda,\lambda+1$ の冪だけであること、代入後の全残余式の一致を確認した。
従って自由係数の特殊値を一般点の除算で落としていない。

$h=x_{15},s=16h^2-1$ と置く。
$(x_2,x_5)=(0,0)$ の全残余は、単元倍を除き
\[
 hx_9=0,\quad sx_6=sx_7=0,\quad
 (\lambda-1)s=0,\quad s\{(\lambda-1)h+x_8\}=0.
\]
$s\ne0$ なら $\lambda=1,x_6=x_7=x_8=0,hx_9=0$。
$s=0$ なら $h=\pm1/4,x_9=0$ であり、
残る三自由係数を保持した二境界族の三次商次元は22となる。

$(0,1)$ では、残余に $hx_6,hx_{17},sx_6,sx_{17}$ があるため
$x_6=x_{17}=0$、残る条件は $h(\lambda-1)s=0$ である。
$h=\pm1/4$ の二境界族は三次商次元21で排除される。

$(1,0)$ では $s\ne0$ の全局所化が
$\lambda=1,h=x_6=x_7=x_8=0$ を強制し、$x_9$ を自由に残す。
$h=\pm1/4$ では全残余が
\[
 x_9=x_6^2,\quad x_6(\lambda h+x_8)=0,\quad
 (\lambda h+x_8)((\lambda+1)h+x_8)=0
\]
となる。従って $x_8=-\lambda h$ と
$x_8=-(\lambda+1)h,x_6=0$ の二枝で全てを覆う。
両符号を合わせた四境界族の三次商次元は、それぞれ21と22である。

$(1,1)$ では $h\ne0$ の全飽和idealが単位idealであり、
$h=0$ とすると $x_6=x_7=0$ を得る。
以上の飽和では $\lambda$ を多項式変数として残し、
$\lambda(\lambda+1)$ と各枝で指定した非零係数だけを逆元化する。
当repoの独立検証器も、四つの主飽和idealを元残余式から再計算した。

八つの境界族の商次元は元六関係から独立に計算した。
元の分母は $\lambda,\lambda+1$ の冪だけであり、
全語関係行列のrank特殊化により、上の下界は全許容閉値へ及ぶ。
従って各族の自由係数が消える点も別途捨てていない。

残る全係数は、$c,e,f,r$ を用いて
\[
\begin{gathered}
 x_1=1,\ x_2=e,\ x_5=f,\ x_9=r,\ x_{11}=-c,\ x_{15}=c,\\
 x_{19}=\lambda/16-c^2,\quad x_{20}=-e/16,\quad g=1/8,\\
 (\lambda-1)c=ce=cr=0,\qquad16fr=\lambda^2-1,\qquad c\ne\pm1/4
\end{gathered}
\]
であり、他の係数は零である。
元の全Schur式への逆代入も、四制約idealでの剰余零として確認した。
$c\ne0$ なら $\lambda=1,e=r=0$ であり、$f=0$ 又は $f=1$ へ移して
$L_c,M_c$ になる。
$c=0,f\ne0$ なら $f=1,e=0$ 又は1で $F^\varepsilon_\lambda$。
$c=f=0$ なら $\lambda=1$ であり、$r\ne0$ を1に移せば $G^\varepsilon$、
$r=0$ なら $L_0$ 又は $K$ となる。
これらの規格化は独立な対角比と根の選択だけを用い、全許容値を覆う。
元16語上の関係は既存FN2表示と文字どおり一致するので、
命題\ref{prop:cubic-fn-lm}と命題\ref{prop:cubic-fn-remaining}により全てH3でない。
\end{proof}

独立検証器は\\
\texttt{verify\_cubic\_CY\_four\_point\_first\_cover.py}。
固定出典の全controlも隔離した出力先で再実行しPASSを確認したが、
AS正則性の十分性や全自己同型の結果は本排除に用いない。
最小固定入力には必要な消去・残余・候補係数だけを保持し、
独立商計算では各次数の全元関係挿入の座標像が零となることも検証する。
$x_0\ne0$ の共有chart、出次数2・部分活性・全共有支持の必要係数被覆は未完成である。
